\begin{theorem}[单 Fourier 模约束下的反向相对熵精确极小与唯一极值族]\label{thm:xi-reversekl-single-frequency-exact-minimizer}
\leanverified{paper\_xi\_reversekl\_single\_frequency\_exact\_minimizer}
固定整数 \(n\ge 1\) 与 \(0<r<1\)。
对任意 \(c\in[0,1]\)，考虑约束集合
\[
\mathcal C_{c,n}:=\{\mu\in\mathrm{Prob}(\TT):\ |\widehat\mu(n)|=c\}.
\]
则有精确极小值
\begin{equation}\label{eq:xi-reversekl-single-frequency-exact-min-value}
\boxed{\;
\inf_{\mu\in\mathcal C_{c,n}}\mathcal H_r(\mu)=-\log\bigl(1-r^{2n}c^2\bigr).\;}
\end{equation}
并且存在如下显式极值族达到该下确界。
对任意 \(\theta_0\in\TT\) 定义概率测度 \(\nu_{c,n,\theta_0}\)：
\[
\dd\nu_{c,n,\theta_0}(\theta)=P_c\!\bigl(n(\theta-\theta_0)\bigr)\,\dd m(\theta)\qquad(0\le c<1),
\qquad
\nu_{1,n,\theta_0}:=\frac1n\sum_{j=0}^{n-1}\delta_{\theta_0+2\pi j/n}.
\]
则 \(\nu_{c,n,\theta_0}\in\mathcal C_{c,n}\) 且
\begin{equation}\label{eq:xi-reversekl-single-frequency-exact-min-attained}
\mathcal H_r(\nu_{c,n,\theta_0})=-\log\bigl(1-r^{2n}c^2\bigr).
\end{equation}
此外，当 \(0<c<1\) 时，\(\mathcal C_{c,n}\) 上的极小化解除旋转参数 \(\theta_0\) 外唯一：若 \(\mu\in\mathcal C_{c,n}\) 取到 \eqref{eq:xi-reversekl-single-frequency-exact-min-value} 的下确界，则存在 \(\theta_0\in\TT\) 使 \(\mu=\nu_{c,n,\theta_0}\)。
\end{theorem}

\begin{proof}
由推论 \ref{cor:xi-reversekl-single-frequency-closed-lower}，
\(
\mathcal H_r(\mu)\ge -\log(1-r^{2n}|\widehat\mu(n)|^2)
\)，
对 \(\mathcal C_{c,n}\) 取下确界得到
\(
\inf_{\mu\in\mathcal C_{c,n}}\mathcal H_r(\mu)\ge -\log(1-r^{2n}c^2)
\)。

为证明可达性，先计算 \(\nu_{c,n,\theta_0}\) 的 Fourier 系数。
当 \(0\le c<1\) 时，由 \(P_c(\theta)=\sum_{k\in\ZZ}c^{|k|}e^{ik\theta}\) 得
\[
\widehat{\nu_{c,n,\theta_0}}(k)
=\int_{\TT}e^{-ik\theta}P_c\!\bigl(n(\theta-\theta_0)\bigr)\,\dd m(\theta)
=
\begin{cases}
c^{|m|}e^{-inm\theta_0},& k=nm,\ m\in\ZZ,\\
0,& n\nmid k,
\end{cases}
\]
从而 \(|\widehat{\nu_{c,n,\theta_0}}(n)|=c\)。
当 \(c=1\) 时，离散测度 \(\nu_{1,n,\theta_0}\) 满足 \(|\widehat{\nu_{1,n,\theta_0}}(n)|=1\) 且 \(\widehat{\nu_{1,n,\theta_0}}(k)=0\)（\(n\nmid k\)），故亦在 \(\mathcal C_{c,n}\) 内。

再计算 Poisson 平滑密度 \(p_r=P_r*\nu_{c,n,\theta_0}\)。
由 \(\widehat{p_r}(k)=r^{|k|}\widehat{\nu_{c,n,\theta_0}}(k)\) 得
\[
\widehat{p_r}(nm)=(r^n c)^{|m|}e^{-inm\theta_0},\qquad \widehat{p_r}(k)=0\ (n\nmid k),
\]
故
\[
p_r(\theta)=P_{r^n c}\!\bigl(n(\theta-\theta_0)\bigr).
\]
于是
\[
\mathcal H_r(\nu_{c,n,\theta_0})
=\int_{\TT}-\log P_{r^n c}\!\bigl(n(\theta-\theta_0)\bigr)\,\dd m(\theta)
=\int_{\TT}-\log P_{r^n c}(\theta)\,\dd m(\theta),
\]
其中最后一步由 \(\theta\mapsto n\theta\) 对 \(m\) 的测度保持性得到。
由 \(P_\rho(\theta)=\frac{1-\rho^2}{|1-\rho e^{i\theta}|^2}\) 与 Jensen 平均恒等式（见本段前述证书推导中的引用），得
\(
\int_{\TT}-\log P_\rho\,\dd m=-\log(1-\rho^2)
\)。
取 \(\rho=r^n c\) 即得 \eqref{eq:xi-reversekl-single-frequency-exact-min-attained}，从而 \eqref{eq:xi-reversekl-single-frequency-exact-min-value} 成立。

最后证明唯一性（\(0<c<1\)）。
设 \(\mu\in\mathcal C_{c,n}\) 取到 \eqref{eq:xi-reversekl-single-frequency-exact-min-value} 的下确界。
则 \(\mu\) 在推论 \ref{cor:xi-reversekl-single-frequency-closed-lower} 的证明链中必须处处取等。
回顾该证明：取 \(Q(z)=1-a z^n\) 并代入定理 \ref{thm:xi-reversekl-exterior-analytic-witness-log} 后得到下界；
其等号条件由定理 \ref{thm:xi-reversekl-dv-poisson-variational} 的最优性刻画给出，必存在常数 \(K>0\) 与某个 \(\theta_0\in\TT\) 使
\[
p_r(\theta)=\frac{K}{|1-ae^{in(\theta-\theta_0)}|^2}.
\]
由归一化 \(\int_{\TT}p_r\,\dd m=1\) 可得 \(K=1-a^2\)，即
\(
p_r(\theta)=P_a(n(\theta-\theta_0))
\)。
另一方面，\(|\widehat\mu(n)|=c\) 且等号必须在 \(a\) 的一维优化处取得，迫使 \(a=r^n c\)。
因此
\(
p_r(\theta)=P_{r^n c}(n(\theta-\theta_0))
\)。
比较 Fourier 系数得到 \(\widehat p_r(k)=0\) 对 \(n\nmid k\)，从而 \(\widehat\mu(k)=0\) 对 \(n\nmid k\)，且对 \(m\in\ZZ\) 有
\(
\widehat\mu(nm)=c^{|m|}e^{-inm\theta_0}
\)。
在紧交换群 \(\TT\) 上，Fourier--Stieltjes 系数唯一确定概率测度，故 \(\mu=\nu_{c,n,\theta_0}\)。证毕。
\end{proof}

\begin{theorem}[等号刚性：单频封口等价于 \(n\)-折叠 Poisson 谱型]\label{thm:xi-reversekl-single-frequency-rigidity-equivalences}
\leanverified{paper\_xi\_reversekl\_single\_frequency\_rigidity\_equivalences}
固定整数 \(n\ge 1\) 与 \(0<r<1\)。
对任意 \(\mu\in\mathrm{Prob}(\TT)\) 令 \(c:=|\widehat\mu(n)|\in(0,1)\)。
则以下命题等价：
\begin{enumerate}
  \item \(\mathcal H_r(\mu)=-\log(1-r^{2n}c^2)\)；
  \item 存在 \(\theta_0\in\TT\) 使 \(p_r(\theta)=(P_r*\mu)(\theta)=P_{r^n c}\!\bigl(n(\theta-\theta_0)\bigr)\)；
  \item \(\mu=\nu_{c,n,\theta_0}\)（定理 \ref{thm:xi-reversekl-single-frequency-exact-minimizer} 中定义）。
\end{enumerate}
\end{theorem}

\begin{proof}
(3)\(\Rightarrow\)(2) 已在定理 \ref{thm:xi-reversekl-single-frequency-exact-minimizer} 的计算中得到。
(2)\(\Rightarrow\)(1) 由 \(\int_{\TT}-\log P_\rho\,\dd m=-\log(1-\rho^2)\) 直接推出。
(1)\(\Rightarrow\)(3) 为定理 \ref{thm:xi-reversekl-single-frequency-exact-minimizer} 的唯一性部分。证毕。
\end{proof}

\begin{theorem}[熵预算下的可达单频幅度上确界与两相实现]\label{thm:xi-reversekl-single-frequency-entropy-budget-sup}
\leanverified{paper\_xi\_reversekl\_single\_frequency\_entropy\_budget\_sup}
固定整数 \(n\ge 1\)、\(0<r<1\)，并令 \(h\in[0,\infty)\)。
定义熵预算集合
\[
\mathcal F_h:=\{\mu\in\mathrm{Prob}(\TT):\ \mathcal H_r(\mu)\le h\}.
\]
则有精确上确界公式
\[
\boxed{\;
\sup_{\mu\in\mathcal F_h}|\widehat\mu(n)|
=\min\Bigl\{1,\ r^{-n}\sqrt{1-e^{-h}}\Bigr\}. \;}
\]
并且极值实现具有分岔结构：
\begin{enumerate}
  \item 若 \(0\le h\le -\log(1-r^{2n})\)，则上确界由 \(\nu_{c,n,\theta_0}\)（除旋转外唯一）达到，其中
  \[
  c=r^{-n}\sqrt{1-e^{-h}}\in[0,1],
  \qquad
  \mathcal H_r(\nu_{c,n,\theta_0})=h;
  \]
  \item 若 \(h\ge -\log(1-r^{2n})\)，则上确界为 \(1\)，由离散测度 \(\nu_{1,n,\theta_0}\) 达到，且
  \(
  \mathcal H_r(\nu_{1,n,\theta_0})=-\log(1-r^{2n})\le h
  \)。
\end{enumerate}
\end{theorem}

\begin{proof}
由推论 \ref{cor:xi-reversekl-single-frequency-closed-lower}，若 \(\mathcal H_r(\mu)\le h\)，则
\[
-\log\bigl(1-r^{2n}|\widehat\mu(n)|^2\bigr)\le h
\quad\Longleftrightarrow\quad
|\widehat\mu(n)|^2\le r^{-2n}(1-e^{-h}).
\]
再与平凡界 \(|\widehat\mu(n)|\le 1\) 合并得到上确界公式。
可达性由定理 \ref{thm:xi-reversekl-single-frequency-exact-minimizer}：取相应 \(c\) 与任意 \(\theta_0\) 即得；当 \(c>1\) 时改取 \(c=1\) 的离散测度。证毕。
\end{proof}

\begin{theorem}[单频线性化的 Legendre--Fenchel 对偶闭式与极值测度]\label{thm:xi-reversekl-single-frequency-legendre-fenchel-closed}
\leanverified{paper\_xi\_reversekl\_single\_frequency\_legendre\_fenchel\_closed}
固定整数 \(n\ge 1\) 与 \(0<r<1\)，并令 \(\alpha:=r^{2n}\in(0,1)\)。
定义对偶泛函
\[
\Psi_{r,n}(\lambda)
:=\sup_{\mu\in\mathrm{Prob}(\TT)}
\Bigl\{\Re\bigl(\lambda\,\overline{\widehat\mu(n)}\bigr)-\mathcal H_r(\mu)\Bigr\},
\qquad \lambda\in\CC.
\]
记 \(t:=|\lambda|\ge 0\) 与 \(u(t):=\sqrt{1+t^2/\alpha}\)。
则
\[
\boxed{
\Psi_{r,n}(\lambda)=
\begin{cases}
u(t)-1+\log\!\Bigl(\dfrac{2}{u(t)+1}\Bigr),& 0\le t\le \dfrac{2\alpha}{1-\alpha},\\[1.1em]
t+\log(1-\alpha),& t\ge \dfrac{2\alpha}{1-\alpha}.
\end{cases}}
\]
并且极值测度可显式给出：
\begin{enumerate}
  \item 若 \(0<t<\dfrac{2\alpha}{1-\alpha}\)，则最优 \(\mu\) 除旋转外唯一，且可取 \(\mu=\nu_{c^\star,n,\theta_0}\)，其中 \(\theta_0\) 满足 \(e^{-in\theta_0}=\lambda/|\lambda|\)，并且
  \[
  c^\star
  =\frac{\sqrt{t^2+\alpha}-\sqrt{\alpha}}{t\sqrt{\alpha}}
  =\frac{\sqrt{t^2+r^{2n}}-r^{n}}{t\,r^{n}}
  \in(0,1).
  \]
  \item 若 \(t>\dfrac{2\alpha}{1-\alpha}\)，则最优由 \(\mu=\nu_{1,n,\theta_0}\) 达到（\(\theta_0\) 同上）。
\end{enumerate}
\end{theorem}

\begin{proof}
由推论 \ref{cor:xi-reversekl-single-frequency-closed-lower}，
\[
\Re\bigl(\lambda\,\overline{\widehat\mu(n)}\bigr)-\mathcal H_r(\mu)
\le t\,|\widehat\mu(n)|+\log\bigl(1-\alpha|\widehat\mu(n)|^2\bigr).
\]
记 \(c:=|\widehat\mu(n)|\in[0,1]\)，则右端不超过
\[
\sup_{0\le c\le 1}\Bigl\{t c+\log(1-\alpha c^2)\Bigr\}.
\]
另一方面，由定理 \ref{thm:xi-reversekl-single-frequency-exact-minimizer}，对任意 \(c\in[0,1]\) 可取 \(\mu=\nu_{c,n,\theta_0}\)（令 \(e^{-in\theta_0}=\lambda/|\lambda|\)）使
\(
\Re(\lambda\overline{\widehat\mu(n)})=t c
\)
且
\(
\mathcal H_r(\mu)=-\log(1-\alpha c^2)
\)。
故
\[
\Psi_{r,n}(\lambda)=\sup_{0\le c\le 1}\Bigl\{t c+\log(1-\alpha c^2)\Bigr\}.
\]
令 \(\phi(c):=t c+\log(1-\alpha c^2)\)。
若极值点在 \((0,1)\)，则 \(\phi'(c)=0\) 等价于
\(
t=\frac{2\alpha c}{1-\alpha c^2}
\)，
其在 \((0,1)\) 有解当且仅当 \(t\le \frac{2\alpha}{1-\alpha}\)。
对应解为所示 \(c^\star\)，代回并令 \(u(t)=\sqrt{1+t^2/\alpha}\) 即得第一分支闭式
\(
\phi(c^\star)=u(t)-1+\log\frac{2}{u(t)+1}
\)。
当 \(t\ge \frac{2\alpha}{1-\alpha}\) 时，\(\phi\) 在 \([0,1]\) 的最大值取于端点 \(c=1\)，从而得到第二分支 \(t+\log(1-\alpha)\)。
极值测度的分类由上述构造与定理 \ref{thm:xi-reversekl-single-frequency-exact-minimizer} 的唯一性给出。证毕。
\end{proof}

\leanverified{paper\_xi\_reversekl\_single\_frequency\_sharp\_fenchel\_inequality}
\begin{corollary}[最锋利的单频 Fenchel 不等式与等号测度分类]\label{cor:xi-reversekl-single-frequency-sharp-fenchel-inequality}
在定理 \ref{thm:xi-reversekl-single-frequency-legendre-fenchel-closed} 的设定下，对任意 \(\mu\in\mathrm{Prob}(\TT)\) 与任意 \(\lambda\in\CC\) 成立
\[
\boxed{\;
\Re\bigl(\lambda\,\overline{\widehat\mu(n)}\bigr)\ \le\ \mathcal H_r(\mu)+\Psi_{r,n}(\lambda).\;}
\]
并且当 \(\lambda\neq 0\) 时，等号成立当且仅当 \(\mu\) 为定理 \ref{thm:xi-reversekl-single-frequency-legendre-fenchel-closed} 中给出的相应极值测度（按 \(t=|\lambda|\) 所属分支区分），且旋转参数 \(\theta_0\) 满足 \(e^{-in\theta_0}=\lambda/|\lambda|\)。
\end{corollary}

\begin{proof}
该不等式为 \(\Psi_{r,n}\) 的 Legendre--Fenchel 定义的直接推论；等号分类由定理 \ref{thm:xi-reversekl-single-frequency-legendre-fenchel-closed} 的极值刻画给出。证毕。
\end{proof}
